\section{Sujet n\textsuperscript{o}\,6}
\noindent\begin{minipage}{0.5\linewidth}\footnotesize Office du Baccalauréat du Cameroun\end{minipage}%
\hfill\begin{minipage}{0.45\linewidth}\footnotesize\raggedleft Examen : \textbf{Baccalauréat} · Série : \textbf{C/E}\\
Épreuve : \textbf{Mathématiques} · Durée : \textbf{4\,h} · Coef. : \textbf{7}\end{minipage}
\par\smallskip{\color{onbo}\rule{\linewidth}{1pt}}

\subsection*{Partie A — Évaluation des ressources \hfill (15 points)}

\noindent\textbf{Exercice 1.} \hfill\textit{(3 points)}\\
Un technicien mesure la production journalière (en \si{\kilo\watt\hour}) d'un petit champ
de panneaux solaires pendant $10$ jours. Il obtient la série suivante :
\[ 18,\ 20,\ 22,\ 19,\ 24,\ 21,\ 23,\ 20,\ 22,\ 21. \]
\begin{enumerate}[label=\arabic*)]
\item Calculer la moyenne $\bar x$ et la variance $V$ de cette série. \hfill\textit{(1,5 pt)}
\item En déduire l'écart-type $\sigma$ (valeur arrondie à $10^{-2}$). \hfill\textit{(0,75 pt)}
\item Déterminer la médiane de la série. \hfill\textit{(0,75 pt)}
\end{enumerate}

\smallskip
\noindent\textbf{Exercice 2.} \hfill\textit{(3 points)}\\
On installe deux types de panneaux : des panneaux de type $A$ et des panneaux de type $B$.
Un panneau $A$ produit \SI{300}{\watt} et coûte \num{45000}~F\,CFA ; un panneau $B$ produit
\SI{400}{\watt} et coûte \num{70000}~F\,CFA. Un client achète $x$ panneaux $A$ et $y$ panneaux $B$.
Il obtient une puissance totale de \SI{5400}{\watt} pour un coût total de \num{870000}~F\,CFA.
\begin{enumerate}[label=\arabic*)]
\item Traduire les données par un système de deux équations à deux inconnues $x$ et $y$. \hfill\textit{(1 pt)}
\item Écrire ce système sous la forme matricielle $M\!\begin{psmallmatrix}x\\y\end{psmallmatrix}=\begin{psmallmatrix}p\\q\end{psmallmatrix}$,
calculer $\det M$ et en déduire que $M$ est inversible. \hfill\textit{(1 pt)}
\item Résoudre le système ; donner le nombre de panneaux de chaque type. \hfill\textit{(1 pt)}
\end{enumerate}

\smallskip
\noindent\textbf{Exercice 3.} \hfill\textit{(4 points)}\\
Soit $f$ définie sur $]0,+\infty[$ par $f(x)=x-\ln x$, de courbe $(\mathcal{C})$.
\begin{enumerate}[label=\arabic*)]
\item Déterminer les limites de $f$ en $0^+$ et en $+\infty$. \hfill\textit{(1 pt)}
\item Étudier les variations de $f$ et dresser son tableau de variations ; en déduire que
$f(x)\ge 1$ pour tout $x>0$. \hfill\textit{(1,5 pt)}
\item À l'aide d'une intégration par parties, calculer
$\displaystyle I=\int_1^{\mathrm e}\ln x\,\dd x$, puis en déduire
$\displaystyle J=\int_1^{\mathrm e}f(x)\,\dd x$. \hfill\textit{(1,5 pt)}
\end{enumerate}

\smallskip
\noindent\textbf{Exercice 4.} \hfill\textit{(5 points)}\\
L'espace est muni d'un repère orthonormé direct $(O\,;\,\vec\imath,\vec\jmath,\vec k)$. Une
toiture portant des panneaux est modélisée par le plan $(P)$ passant par les points
$A(2\,;\,0\,;\,0)$, $B(0\,;\,3\,;\,0)$ et $C(0\,;\,0\,;\,6)$.
\begin{enumerate}[label=\arabic*)]
\item Déterminer les coordonnées des vecteurs $\vec{AB}$ et $\vec{AC}$, puis montrer que
$\vec n(9\,;\,6\,;\,3)$ est un vecteur normal au plan $(P)$. \hfill\textit{(1,5 pt)}
\item En déduire une équation cartésienne du plan $(P)$. \hfill\textit{(1 pt)}
\item Le soleil est repéré par la direction du vecteur $\vec u(0\,;\,0\,;\,1)$ (verticale).
Calculer une mesure de l'angle $\theta$ entre $\vec n$ et $\vec u$, arrondie au degré. \hfill\textit{(1,5 pt)}
\item Calculer la distance du point $O$ au plan $(P)$. \hfill\textit{(1 pt)}
\end{enumerate}

\subsection*{Partie B — Évaluation des compétences \hfill (5 points)}

\noindent\textit{Madame Ngono fait installer un système d'énergie solaire dans sa maison à
Ebolowa. Elle te sollicite, en tant qu'élève de Terminale~C, pour trois décisions
techniques. Un panneau délivre une puissance de \SI{350}{\watt} en plein soleil ; à cause
de la poussière, son rendement diminue de \SI{2}{\percent} chaque mois sans entretien. Par
ailleurs, elle dispose d'un budget de \num{1200000}~F\,CFA pour acheter des panneaux à
\num{75000}~F\,CFA l'unité et une batterie à \num{300000}~F\,CFA. Enfin, sur un lot de
panneaux livrés, \SI{4}{\percent} présentent un micro-défaut ; un contrôleur en teste $5$ au hasard.}

\begin{enumerate}[label=\textbf{Tâche \arabic*.},leftmargin=2.4em]
\item Déterminer le nombre de mois au bout duquel la puissance d'un panneau non entretenu
sera tombée en dessous de \SI{300}{\watt}. \hfill\textit{(1,5 pt)}
\item Sachant qu'elle achète obligatoirement une batterie, déterminer le nombre maximal de
panneaux qu'elle peut acheter avec son budget. \hfill\textit{(1,5 pt)}
\item Déterminer la probabilité qu'\emph{au moins un} des $5$ panneaux testés soit défectueux. \hfill\textit{(1,5 pt)}
\end{enumerate}
\par\smallskip{\footnotesize\itshape Présentation générale : 0,5 point.}

% ════════════════════════════════════════════════════
\corrige{du sujet 6}

\noindent\textbf{Partie A — Exercice 1.}
1) Somme $=18{+}20{+}22{+}19{+}24{+}21{+}23{+}20{+}22{+}21=210$, donc $\bar x=\frac{210}{10}=21$.
Variance : $V=\frac1{10}\sum(x_i-\bar x)^2$. Les écarts au carré valent
$9,1,1,4,9,0,4,1,1,0$, de somme $30$ ; d'où $V=\frac{30}{10}=3$.
2) $\sigma=\sqrt3\approx1{,}73$.
3) Série ordonnée : $18,19,20,20,21,21,22,22,23,24$. Médiane $=\frac{21+21}{2}=21$.

\noindent\textbf{Exercice 2.}
1) Puissance : $300x+400y=5400$, soit $3x+4y=54$. Coût : $45000x+70000y=870000$,
soit (en divisant par $5000$) $9x+14y=174$.
2) $M=\begin{psmallmatrix}3&4\\9&14\end{psmallmatrix}$, $\begin{psmallmatrix}p\\q\end{psmallmatrix}=\begin{psmallmatrix}54\\174\end{psmallmatrix}$.
$\det M=3\times14-4\times9=42-36=6\neq0$ : $M$ est inversible.
3) $M^{-1}=\frac16\begin{psmallmatrix}14&-4\\-9&3\end{psmallmatrix}$. Donc
$x=\frac16(14\times54-4\times174)=\frac16(756-696)=\frac{60}{6}=10$ et
$y=\frac16(-9\times54+3\times174)=\frac16(-486+522)=\frac{36}{6}=6$.
Le client achète \textbf{$10$ panneaux $A$} et \textbf{$6$ panneaux $B$}
(vérification : $300\cdot10+400\cdot6=5400$ ; $45000\cdot10+70000\cdot6=870000$).

\noindent\textbf{Exercice 3.}
1) En $0^+$, $\ln x\to-\infty$ donc $f(x)=x-\ln x\to+\infty$. En $+\infty$,
$f(x)=x(1-\frac{\ln x}{x})\to+\infty$ (car $\frac{\ln x}{x}\to0$).
2) $f'(x)=1-\frac1x=\frac{x-1}{x}$ : $f'<0$ sur $]0,1[$, $f'>0$ sur $]1,+\infty[$ ; minimum en $x=1$
avec $f(1)=1-\ln1=1$. Donc $f(x)\ge f(1)=1$ pour tout $x>0$.
3) IPP ($u=\ln x$, $v'=1$, $u'=\frac1x$, $v=x$) :
$I=[x\ln x]_1^{\mathrm e}-\int_1^{\mathrm e}1\,\dd x=(\mathrm e\cdot1-0)-(\mathrm e-1)=1$.
Puis $J=\int_1^{\mathrm e}x\,\dd x-\int_1^{\mathrm e}\ln x\,\dd x=[\frac{x^2}{2}]_1^{\mathrm e}-I
=\frac{\mathrm e^2-1}{2}-1=\frac{\mathrm e^2-3}{2}\approx2{,}69$.

\noindent\textbf{Exercice 4.}
1) $\vec{AB}=B-A=(-2\,;\,3\,;\,0)$, $\vec{AC}=C-A=(-2\,;\,0\,;\,6)$.
$\vec n\cdot\vec{AB}=9(-2)+6\cdot3+3\cdot0=-18+18=0$ et
$\vec n\cdot\vec{AC}=9(-2)+6\cdot0+3\cdot6=-18+18=0$ : $\vec n$ est orthogonal à deux vecteurs
non colinéaires de $(P)$, c'est donc un vecteur normal.
2) $(P)\!:\ 9x+6y+3z+d=0$ ; avec $A(2,0,0)$ : $18+d=0$, $d=-18$. Soit $9x+6y+3z-18=0$,
ou $3x+2y+z-6=0$.
3) On simplifie $\vec n=(9,6,3)$ en $\vec n'=(3,2,1)$ (colinéaire). Alors
$\cos\theta=\dfrac{|\vec n'\cdot\vec u|}{\|\vec n'\|\,\|\vec u\|}=\dfrac{|1|}{\sqrt{9+4+1}\cdot1}=\dfrac{1}{\sqrt{14}}\approx0{,}267$,
d'où $\theta\approx75^{\circ}$.
4) $d(O,(P))=\dfrac{|3\cdot0+2\cdot0+0-6|}{\sqrt{3^2+2^2+1^2}}=\dfrac{6}{\sqrt{14}}\approx1{,}60$.

\smallskip
\noindent\textbf{Partie B — Situation.}
\emph{Tâche 1.} Puissance après $n$ mois : $P_n=350\times0{,}98^{\,n}$. On cherche le plus petit
$n$ tel que $P_n<300$, soit $0{,}98^{\,n}<\frac{300}{350}=\frac67\approx0{,}857$. Donc
$n>\dfrac{\ln(6/7)}{\ln0{,}98}\approx\dfrac{-0{,}1542}{-0{,}0202}\approx7{,}6$ : au bout de
\textbf{$8$ mois}.

\emph{Tâche 2.} Avec une batterie obligatoire, le budget restant pour les panneaux est
$1\,200\,000-300\,000=900\,000$. Nombre de panneaux $\le\frac{900\,000}{75\,000}=12$. Elle peut
acheter au maximum \textbf{$12$ panneaux}.

\emph{Tâche 3.} $P(\text{au moins un défectueux})=1-(0{,}96)^5\approx1-0{,}8154=\mathbf{0{,}185}$,
soit environ \SI{18,5}{\percent}.
